\section{Questions raised in review, and what the evidence says}\label{sec:notes}

Review of the previous version raised ten questions. Each is stated here with the evidence found, a verdict, and the
change it made to the paper. New analyses are experiments E18--E21 (Section~\ref{sec:experiments}); new formal results are
in Section~\ref{sec:modelix}.

\subsubsection{1. Does ``female choice'' assume consent? Was rape common}
\emph{Evidence.} Table~\ref{tab:cultures} gives every relevant code in the Standard Cross-Cultural Sample. A daughter
chooses alone in \CcAloneW\% of societies; her consent is unnecessary in two-thirds of those with information; rape is
reported in about half of the societies coded for it by one team and rated common in \CcRapeCom{} of \CcRapeN{} by another;
captives, wife-beating and unequal access to divorce are all recorded. Among the \CcRapeParN{} societies coded for both,
rape is rated common in \CcRapeParCom{} where parents choose a daughter's husband and in \CcRapeOwnCom{} of \CcRapeOwnN{}
where she chooses; the numbers are too small to bear weight. Low rates of extra-pair paternity in historical
populations~\cite{larmuseau2016} bound conceptions outside recognized unions, not coercion within them. Male aggression
against women has a comparative and evolutionary literature of its own~\cite{smuts1992}, and raiding for women is described
ethnographically~\cite{chagnon1988}. \emph{Verdict.} The objection is right about the word and does not touch the
mechanism. No rate of coercion can be given for the Pleistocene. \emph{Change.} The paper no longer writes of choice by
women where it means decision by whoever decides; Eq.~\eqref{eq:coercion} states what coercion does to the pathway;
Section~\ref{sec:reschoice} reports the record; objection O45 and prediction P30 follow.

\begin{table*}[t]
\caption{Who decides a marriage, and how much constraint and coercion is coded, in the Standard Cross-Cultural Sample and in
Binford's hunter-gatherer compilation (D-PLACE release~\cite{kirby2016}). Counts and shares of societies with a code; no
significance is claimed across societies, which share ancestry and neighbors.}\label{tab:cultures}
\scriptsize
\begin{tabularx}{\textwidth}{@{}P{0.24\textwidth}P{0.10\textwidth}rL@{}}
\toprule
\tabinput{generated/tab_cultures.tex}
\bottomrule
\end{tabularx}
\end{table*}

\subsubsection{2. Did only half of all men reproduce? Is there twice as much female DNA}
\emph{Evidence.} The claim descends from the twofold deeper coalescence of mitochondrial than Y-chromosome
lineages~\cite{wilder2004} and was popularized as ``80\% of women and 40\% of men''~\cite{baumeister2010}. Later work with
complete sequences confirms larger differences among populations for the Y chromosome, with regional variation and smaller
effects than first reported~\cite{lippold2014}; X-linked diversity is higher than expected relative to
autosomes~\cite{hammer2008}; a recombination-based estimate puts the breeding ratio nearer one than two~\cite{labuda2010}; and a
bottleneck in Y lineages within the last ten thousand years~\cite{karmin2015} is explained by competition among patrilineal
kin groups without any change in how many men fathered children~\cite{zeng2018}. Across societies with data, the variance
in men's reproductive success exceeds women's by amounts that depend on the marriage system~\cite{brown2009,betzig2012}.
\emph{Verdict.} Half right. Men's reproduction has been the more unequal; the headcount is not recoverable from the ratio
(Proposition~\ref{prop:ancestors}); and autosomal DNA comes equally from mothers and fathers in every generation, so the
second claim is false as stated. \emph{Change.} Proposition~\ref{prop:ancestors}, Fig.~\ref{fig:choice}D, estimates for
the United States, and the link to the floor (Corollary~\ref{cor:floor}).

\subsubsection{3. Did \textit{Homo floresiensis} live beside australopiths, or alone on an island}
\emph{Evidence.} Neither beside australopiths nor beside anyone. The Liang Bua remains date to about 100,000--60,000 years
ago~\cite{sutikna2016}; similar hominins were on Flores by 0.7 million years ago~\cite{brumm2016,vandenbergh2016}; a second
small-bodied island species is known from Luzon~\cite{detroit2019}. The last australopiths lived more than a million years
earlier, in Africa. The brain of \textit{H.~floresiensis} is of australopith size~\cite{kubo2013}, which is the source of
the confusion and was the point of plotting it. \emph{Verdict.} The earlier wording (``persisted beside large-brained
ones'') invited the misreading. \emph{Change.} Text corrected in Sections~\ref{sec:resfossil} and~\ref{sec:reshominins};
isolation is now part of the explanation offered.

\subsubsection{4. Was attention deficit an advantage for hunter-gatherers? Is there a 20\% figure}
\emph{Evidence.} Section~\ref{sec:reshominins} gives the five findings we could verify~\cite{ding2002,matthews2011,eisenberg2008,barack2024,esteller2020}
and one theory of how a minority strategy could persist~\cite{williams2006}. The Ariaal study compared 65 nomadic and 87
settled men and found body mass index higher in carriers of the \textit{DRD4} 7R allele among nomads and lower among the
settled, mainly through fat-free mass; it reports no percentage advantage in efficiency. \emph{Verdict.} Plausible for
exploration-related traits in mobile settings; unsupported as a claim that a clinical condition was selected for; the 20\%
figure could not be traced. \emph{Change.} Corollary~\ref{cor:narrow} states what the framework can say: a mixed profile is
neutral under broad comparison and a deficit under a narrow test. Prediction P33 tests it.

\subsubsection{5. Brains shrank while cognitive alleles rose: doing more with less}
\emph{Evidence.} Both parts hold with caveats: a decline of more than 5\% in 50,000 years, largely tracking body
size~\cite{stibel2021,ruff1997}, with the timing of any later drop disputed~\cite{desilva2021,villmoare2022}; and a
one-standard-deviation rise over ten thousand years in combinations of alleles that predict cognitive measures
today~\cite{akbari2026}. Selection for prosociality, likened to domestication, has also been proposed for recent human
evolution~\cite{hare2017}. \emph{Verdict.} Compatible, and predicted by Proposition~\ref{prop:tissue} when tissue is costly. The
comparison with efficient machine-learning models is an analogy of form. \emph{Change.} A paragraph in
Section~\ref{sec:reshominins} and a statement in Section~\ref{sec:ai}.

\subsubsection{6. Was neurodivergence selected for the forest: sharper senses, attention, smell}
\emph{Evidence.} Hunter-gatherers name odors far better than neighboring horticulturalists who speak a closely related
language~\cite{majid2018}, which attributes the difference to way of life. Common variants associated with autism show
signatures of positive selection, which the authors relate to cognitive ability~\cite{polimanti2017}. We found
no study showing heritable sensory acuity that was favored in foragers and relaxed in farmers. \emph{Verdict.} Not
supported as stated. Diagnostic categories are modern and should not be projected onto past populations. \emph{Change.}
One sentence in Section~\ref{sec:reshominins}.

\subsubsection{7. Do storytellers have more children? What of social networks, fire, and staying up late}
\emph{Evidence.} Yes for Agta storytellers~\cite{smith2017}, for network position among Agta and BaYaka
mothers~\cite{page2017}, and for cooperative ties among the BaYaka~\cite{chaudhary2016}; Hadza networks are structured by
cooperation~\cite{apicella2012}; firelit talk is stories~\cite{wiessner2014,dunbar2014}, and people without electricity
stay awake some hours after dark~\cite{yetish2015}. Status of several kinds goes with men's reproductive success across
33 nonindustrial societies~\cite{vonrueden2016}. We could not identify the book cited to us as \emph{Outbound}. \emph{Verdict.}
Supported, and it enters the model in two places that the cooking hypothesis does not: fire raises the number of occasions
and observers (Eq.~\ref{eq:sb}) and adds speech as a domain of common comparison (Corollary~\ref{cor:generalist}).
\emph{Change.} A paragraph in Section~\ref{sec:resjury}; prediction P31.

\subsubsection{8--9. The Maasai: lion hunting replaced by jumping; safe, repeatable rituals}
\emph{Evidence.} The tournament exists and was created for that purpose~\cite{maasaiolympics,aljazeera2022}; lion killing by
Maasai has several motives, of which the warrior's proof is one~\cite{hazzah2009,goldman2013}; marriage by custom follows
the ceremony that ends warriorhood~\cite{aljazeera2023,spencer1988}. Ritual fights raise men's reproductive success among the
Sereer~\cite{llaurens2009}, sports are overwhelmingly male in 50 societies~\cite{deaner2013}, and games with winners are
nearly universal (\CcGamesAny{} of \CcGamesN{} societies; \CcGamesForAny{} of \CcGamesForN{} foraging societies). We report
only the presence of games~\cite{roberts1959}: the distribution of types in the release looked implausible to us (most
foraging societies coded as having games of strategy alone), and we could not check it against the source. \emph{Verdict.} Supported as a case of substitution.
That jump height decides eligibility to marry is not supported by any source we found; prestige is what is at stake.
\emph{Change.} Corollary~\ref{cor:ritual}, a paragraph in Section~\ref{sec:rescontest}, prediction P32.

\subsubsection{10. Specialization, online life, the loss of community, and falling births}
\emph{Evidence.} Tables~\ref{tab:births} and~\ref{tab:village}. Ideals have not moved; the fall in children per young
adult is accounted for almost wholly by the rise of the never married; social evenings with neighbors and friends, and
weekly worship, have declined; time online has risen; having no partner has become commoner, especially for men. Against
the simple story: seeing friends goes with being single, time on social media is unrelated to being single in 70
countries, and couples who met online married sooner in one panel~\cite{rosenfeld2017}. Completed fertility falls short of
intentions in Europe and the United States generally~\cite{beaujouan2019}, and social networks shape fertility
decisions~\cite{bernardi2014}. \emph{Verdict.} Supported at the level of accounting and trend; the causal step from thin
community to fewer unions is conjecture, formalized as Proposition~\ref{prop:waiting}. On generalists: the model says
division of labor removes the premium on general ability (Proposition~\ref{prop:division}); it does not say generalists
survive better. \emph{Change.} Section~\ref{sec:resvillage}, Fig.~\ref{fig:village}, predictions P34 and P35, objection O49.

\begin{table*}[t]
\caption{Marriage, children and ideals among adults in the General Social Survey. Children and marital status: ages
25--39. Ideal number of children for a family: ages 18--44, answers of ``as many as you want'' excluded. Weighted.}\label{tab:births}
\scriptsize
\begin{tabularx}{\textwidth}{@{}lrCCCCC@{}}
\toprule
\tabinput{generated/tab_births.tex}
\bottomrule
\end{tabularx}
\end{table*}

\begin{table*}[t]
\caption{Social ties and partnership among adults aged 22--34, by survey period (General Social Survey; percentages unless
stated; weighted). Questions on partners began in 1988 and on time online in 2000. Surveys from 2021 were partly
self-administered.}\label{tab:village}
\scriptsize
\begin{tabularx}{\textwidth}{@{}Lrrrrr@{}}
\toprule
\tabinput{generated/tab_village.tex}
\bottomrule
\end{tabularx}
\end{table*}

\begin{table*}[t]
\caption{Stated ideals and partnership across countries. K: 175-country survey~\cite{kowal2025}, countries with at least 50
men and 50 women (150 of each for the partnership row). W: 45-country survey~\cite{walter2020}. Differences $d$ are in pooled
standard deviations. No country is named or ranked.}\label{tab:world}
\scriptsize
\begin{tabularx}{\textwidth}{@{}LcrP{0.34\textwidth}@{}}
\toprule
\tabinput{generated/tab_world.tex}
\bottomrule
\end{tabularx}
\end{table*}

\subsubsection{The conclusions were not strong or important enough}
The earlier version reported what could be derived and measured and stopped. This one states eight conclusions and marks
the status of each (Section~\ref{sec:conclude}). The two with the widest consequences are that the reading of minds is a
property of communities, and that the births now missing in the United States are missing at the step where a couple
forms. Both are stated as strongly as the data allow and no more strongly: the first is measured and derived, the second
is an accounting in one country with a conjectured mechanism and three findings against its simplest form.

\subsubsection{Further open data examined}
Of fifty candidate sources suggested to us, two were added as analyses: the 175-country survey~\cite{kowal2025}, licensed
CC~BY~4.0, and the 45-country survey~\cite{walter2020} (Table~\ref{tab:world}; E20). Published findings on length of
acquaintance~\cite{hunt2015} and on romance that begins as friendship~\cite{stinson2022} were added as evidence. Public-domain
deposits of face-rating data could not be downloaded in our build environment, and panels that require a signed agreement
(on how couples meet, and national fertility surveys) were not used; they are the natural next tests of predictions P28,
P34 and P35.
